\chapter{Implementación y Resultados}
\label{cap:resultados}

Este capítulo constituye el núcleo técnico del proyecto. Se describirá en profundidad la implementación funcional de cada uno de los subsistemas del \textit{Enterprise RAG System}, acompañando las explicaciones con fragmentos de código reales extraídos del repositorio, tablas de configuración y ejemplos prácticos de uso.

% ============================================================
\section{Arquitectura del Agente Enrutador JARVIS}
% ============================================================

El agente \textbf{JARVIS} es la pieza central de orquestación del sistema. Implementado como una clase Python \texttt{Pipeline} compatible con la interfaz de plugins de OpenWebUI, actúa como un enrutador inteligente que analiza cada mensaje del usuario y decide a qué modo de operación derivar la petición.

\subsection{Detección de Intenciones}
El mecanismo de detección de intenciones se implementa mediante un sistema de reglas heurísticas basado en coincidencia de palabras clave y expresiones regulares, evaluadas en un orden estrictamente prioritario. Se definen las siguientes diez categorías de intención:

\begin{table}[H]
\centering
\caption{Modos de operación del agente JARVIS y sus criterios de activación.}
\label{tab:jarvis-modos}
\begin{tabular}{lll}
\toprule
\textbf{Modo} & \textbf{Acción} & \textbf{Ejemplo de activación} \\
\midrule
Visión/OCR     & \texttt{ocr}          & Imagen adjunta en el chat \\
RAG Documental & \texttt{rag}          & ``¿Qué dice el manual de calidad?'' \\
Web Search     & \texttt{web\_search}  & ``Busca en internet sobre Docker'' \\
Web Scraping   & \texttt{web\_scrape}  & Incluir una URL en el mensaje \\
BOE Legal      & \texttt{boe\_search}  & ``Busca en el BOE la ley de datos'' \\
Listar Docs    & \texttt{list\_docs}   & ``/docs'' o ``¿Qué documentos tienes?'' \\
Status         & \texttt{check\_status}& ``¿Cómo va la indexación?'' \\
Ayuda          & \texttt{help}         & ``¿Qué puedes hacer?'' \\
Conversación   & \texttt{chat}         & Cualquier otro mensaje \\
\bottomrule
\end{tabular}
\end{table}

La evaluación sigue un protocolo de prioridades estricto: primero se verifican las condiciones más específicas (como la presencia de archivos adjuntos o imágenes), y se desciende progresivamente hacia las heurísticas más genéricas. Este diseño evita que un modo ambiguo como Web Search capture consultas que deberían derivarse al RAG.

\subsection{Memoria de Sesión por Usuario}
JARVIS mantiene un estado volátil en memoria para cada usuario identificado por su dirección de correo electrónico. Este estado incluye tres estructuras fundamentales:

\begin{itemize}
    \item \texttt{\_user\_image\_memory}: Almacena la última imagen enviada por el usuario en formato Base64, permitiendo preguntas de seguimiento como ``¿Qué pone en la imagen anterior?'' sin necesidad de reenviarla.
    \item \texttt{\_user\_web\_memory}: Retiene el contenido extraído de la última URL analizada, habilitando consultas contextuales como ``Resume la web que acabamos de mirar''.
    \item \texttt{\_user\_pending\_decision}: Gestiona decisiones pendientes cuando el sistema detecta que una URL ya fue indexada previamente y ofrece al usuario la opción de reutilizar el contenido existente o volver a scrapearla.
\end{itemize}

\subsection{Control de Acceso Basado en Grupos}
El pipeline integra un sistema de permisos multi-departamento que segmenta el acceso a las colecciones vectoriales según los grupos de Azure Active Directory del usuario. Al recibir una petición, JARVIS extrae los grupos del objeto \texttt{user\_data} proporcionado por OpenWebUI y los mapea a colecciones de Qdrant mediante un diccionario configurable:

\begin{lstlisting}[language=Python, caption={Configuración del mapeo departamental en JARVIS.}]
DEPARTMENT_MAPPING: Dict[str, str] = {
    "CIVEX2": "documents_CIVEX2",
    "7573b3c1-eeb0-4e3b-8c41-08b749a1dffd": "documents_CIVEX2",
}
DEFAULT_COLLECTION: str = "documents"
\end{lstlisting}

Si el usuario no pertenece a ningún grupo específico, accede únicamente a la colección por defecto. Los usuarios con rol de administrador acceden automáticamente a todas las colecciones. Además, se implementa una caché de permisos con un \textit{Time-To-Live} (TTL) configurable de 300 segundos para minimizar las llamadas repetitivas a la API de Microsoft Graph.

% ============================================================
\section{Pipeline de Embeddings y Búsqueda Semántica}
% ============================================================

\subsection{Modelo de Embeddings}
El sistema utiliza el modelo \texttt{paraphrase-multilingual-MiniLM-L12-v2} de la librería \textit{SentenceTransformers}, que genera vectores densos de 384 dimensiones optimizados para la similitud del coseno multilingüe. Este modelo fue seleccionado por los siguientes criterios técnicos:

\begin{table}[H]
\centering
\caption{Comparativa de modelos de embeddings evaluados.}
\label{tab:embeddings-comparativa}
\begin{tabular}{lccc}
\toprule
\textbf{Modelo} & \textbf{Dimensiones} & \textbf{Multilingüe} & \textbf{Tamaño (MB)} \\
\midrule
all-MiniLM-L6-v2           & 384  & No  & 80  \\
paraphrase-multilingual-MiniLM-L12-v2 & 384 & Sí & 420 \\
all-mpnet-base-v2          & 768  & No  & 420 \\
multilingual-e5-large      & 1024 & Sí  & 2200 \\
\bottomrule
\end{tabular}
\end{table}

La ejecución se fuerza explícitamente en CPU dentro del contenedor del backend para evitar conflictos de versiones CUDA con PaddleOCR, que monopoliza la GPU para tareas de reconocimiento óptico:

\begin{lstlisting}[language=Python, caption={Inicialización del generador de embeddings forzando ejecución en CPU.}]
class EmbeddingGenerator:
    def __init__(self, model_name="paraphrase-multilingual-MiniLM-L12-v2"):
        self.model = SentenceTransformer(model_name, device="cpu")
        self.vector_dim = self.model.get_sentence_embedding_dimension()
\end{lstlisting}

\subsection{Estrategia de Chunking}
Los documentos ingeridos se segmentan en fragmentos (\textit{chunks}) de tamaño configurable, típicamente entre 512 y 1024 tokens, con un solapamiento (\textit{overlap}) del 10--15\% entre fragmentos consecutivos. Este solapamiento es vital para mantener la coherencia semántica en párrafos que cruzan los límites de ventana.

\subsection{Búsqueda Híbrida en Qdrant}
La recuperación de contexto combina dos estrategias complementarias:
\begin{enumerate}
    \item \textbf{Búsqueda densa HNSW:} Emplea los vectores de 384 dimensiones para localizar los fragmentos semánticamente más próximos mediante el índice jerárquico HNSW (\textit{Hierarchical Navigable Small World}).
    \item \textbf{Búsqueda dispersa BM25:} Complementa la anterior con una búsqueda léxica clásica que favorece coincidencias exactas de términos, especialmente útil para códigos de documento, siglas o identificadores.
\end{enumerate}

% ============================================================
\section{Extracción Web (Scraping)}
% ============================================================

\subsection{Motor de Renderizado con Playwright}
La extracción web se implementa a través de un módulo basado en \textit{Playwright}, un framework de automatización de navegadores desarrollado por Microsoft. A diferencia de soluciones simplistas basadas en peticiones HTTP, Playwright instancia un navegador Chromium completo en modo \textit{headless} (sin interfaz gráfica), lo cual permite:

\begin{itemize}
    \item Renderizar páginas JavaScript complejas (React, Vue, Angular).
    \item Ejecutar la carga diferida (\textit{lazy loading}) de contenido.
    \item Resolver redirecciones y consentimientos de cookies.
    \item Capturar el DOM final tras la ejecución de todos los scripts.
\end{itemize}

El scraper incorpora rotación de \textit{User-Agent}, reintentos con retroceso exponencial (\textit{exponential backoff}) y una espera inteligente configurable para aplicaciones de una sola página (SPA). Tras obtener el HTML renderizado, se aplica una cadena de extracción de contenido compuesta por tres estrategias de \textit{fallback}:

\begin{enumerate}
    \item \textbf{Trafilatura:} Algoritmo NLP especializado en extraer contenido principal de páginas web eliminando menús, cabeceras, pies de página y publicidad. Se prioriza por su alta precisión.
    \item \textbf{Readability:} Algoritmo alternativo basado en heurísticas de densidad de texto, comparable al modo lectura de los navegadores modernos.
    \item \textbf{BeautifulSoup:} Extracción cruda de todo el texto visible del DOM como última alternativa si las dos anteriores no producen un resultado superior al umbral mínimo de 200 caracteres.
\end{enumerate}

\subsection{Búsqueda en Internet}
Complementariamente al scraping de URLs específicas, el sistema ofrece un modo de búsqueda general en internet. Mediante la librería \texttt{duckduckgo\_search} y un \textit{fallback} propio de scraping HTML directo sobre el dominio \texttt{html.duckduckgo.com}, se ejecutan consultas en lenguaje natural y se devuelven los resultados más relevantes con sus enlaces directos, evitando redirects intermedios.

% ============================================================
\section{Reconocimiento Óptico de Caracteres (OCR)}
% ============================================================

\subsection{Motor PaddleOCR con Aceleración GPU}
El motor de OCR seleccionado es \textbf{PaddleOCR}, desarrollado por Baidu, que ofrece reconocimiento de texto en más de 80 idiomas con soporte de ejecución acelerada por GPU mediante CUDA. Se ejecuta sobre las tarjetas NVIDIA del servidor utilizando los núcleos CUDA para el procesamiento paralelo de las imágenes.

La inicialización del motor implementa un mecanismo de \textit{fallback} automático: si la GPU no está disponible o falla la carga de los controladores CUDA, se reinicia automáticamente en modo CPU sin intervención manual:

\begin{lstlisting}[language=Python, caption={Inicialización de PaddleOCR con fallback CPU/GPU.}]
try:
    self.ocr = PaddleOCR(
        use_angle_cls=True,
        lang="latin",
        use_gpu=True,
        gpu_mem=500
    )
except Exception as e:
    logger.warning(f"Fallo GPU, reintentando en CPU: {e}")
    self.ocr = PaddleOCR(
        use_angle_cls=True,
        lang="latin",
        use_gpu=False
    )
\end{lstlisting}

\subsection{Flujo de Procesamiento de PDFs Escaneados}
Cuando el backend recibe un archivo PDF, ejecuta un proceso de detección inteligente en dos fases:
\begin{enumerate}
    \item \textbf{Extracción textual directa:} Se intenta obtener texto mediante bibliotecas de parseo de PDF como \texttt{pdfplumber}. Si el volumen de texto extraído es satisfactorio ($> 100$ caracteres por página), se asume que el documento es nativo digital y se procesa directamente.
    \item \textbf{Fallback OCR:} Si la extracción directa produce un resultado insuficiente (indicador de que se trata de un PDF escaneado), cada página se convierte a imagen y se pasa por el motor PaddleOCR, que detecta los cuadros delimitadores de texto, los clasifica angularmente y transcribe su contenido.
\end{enumerate}

% ============================================================
\section{Integración con Microsoft SharePoint}
% ============================================================

El módulo \textbf{Indexer} funciona como un servicio autónomo (\textit{worker}) que ejecuta un ciclo de sincronización periódica cada 5 minutos mediante un cron interno. Su objetivo es mantener actualizada la base de conocimiento del RAG con los documentos corporativos alojados en SharePoint.

\subsection{Autenticación MSAL y Microsoft Graph}
La conexión con SharePoint se realiza mediante el protocolo OAuth 2.0 con flujo de credenciales de cliente (\textit{Client Credentials Flow}), implementado a través de la librería MSAL (\textit{Microsoft Authentication Library}). El servicio obtiene un token de acceso del \textit{tenant} de Azure AD y utiliza la API REST de Microsoft Graph para listar y descargar los ficheros de los sitios configurados.

\subsection{Sincronización Delta Incremental}
Para evitar reprocesar documentos ya ingeridos, el Indexer implementa un mecanismo de sincronización delta: consulta la marca temporal de modificación de cada fichero en SharePoint y la compara con los registros almacenados en PostgreSQL. Solo los ficheros nuevos o modificados desde la última ejecución son descargados, procesados (chunking + embeddings) e indexados en Qdrant.

% ============================================================
\section{Servidor del Boletín Oficial del Estado (BOE)}
% ============================================================

\subsection{Implementación del Protocolo MCP}
La integración con el BOE se implementa utilizando el \textit{Model Context Protocol} (MCP), un estándar abierto propuesto por Anthropic para formalizar la comunicación entre modelos de lenguaje y herramientas externas. Se despliega un servidor MCP local construido con la librería \texttt{fastmcp}.

El servidor expone funciones de búsqueda y consulta que transforman las peticiones en lenguaje natural del usuario en llamadas estructuradas a la API de Datos Abiertos del BOE Español.

\subsection{Resolución Semántica de Leyes}
Para consultas del tipo ``Resume la Ley de Protección de Datos'', el sistema implementa un mecanismo de resolución semántica que mapea nombres coloquiales a identificadores oficiales del BOE:

\begin{lstlisting}[language=Python, caption={Mapeo de nombres comunes a identificadores BOE.}]
LAW_MAPPINGS = {
    "proteccion de datos": "BOE-A-2018-16673",  # LOPDGDD
    "constitucion": "BOE-A-1978-31229",
    "propiedad intelectual": "BOE-A-1996-8930",
    "igualdad": "BOE-A-2007-6115",
}
\end{lstlisting}

El resolvedor elimina palabras vacías (\textit{stop words}) comunes del español (``la'', ``de'', ``del'', ``ley'') y busca coincidencias parciales en el diccionario de mapeos antes de acudir al \textit{endpoint} de búsqueda general del BOE.

% ============================================================
\section{Fine-Tuning con LoRA}
% ============================================================

\subsection{Generación del Dataset de Entrenamiento}
Para adaptar el modelo base a las necesidades específicas del dominio corporativo, se implementó un pipeline de generación automática de datos de entrenamiento. Un script extrae pares pregunta-respuesta directamente de los fragmentos indexados en Qdrant, utilizando el modelo base LLaMA 3.1 como generador sintético de consultas plausibles para cada fragmento.

\subsection{Entrenamiento con Low-Rank Adaptation}
El ajuste fino se realiza mediante la técnica LoRA (\textit{Low-Rank Adaptation}), que inserta matrices de rango reducido en las capas de atención del transformador sin modificar los pesos originales del modelo. Esto permite un entrenamiento eficiente con requisitos de VRAM significativamente inferiores al ajuste fino completo.

El modelo resultante, \texttt{rag-qwen-ft:latest}, se basa en Qwen 2.5 y se registra en LiteLLM como el modelo primario bajo el alias ``JARVIS'' para todas las consultas RAG corporativas.

\begin{table}[H]
\centering
\caption{Modelos de IA desplegados en el sistema y su configuración.}
\label{tab:modelos-ia}
\begin{tabular}{lccl}
\toprule
\textbf{Modelo} & \textbf{Cuantización} & \textbf{VRAM} & \textbf{Uso} \\
\midrule
LLaMA 3.1 8B Instruct & Q8\_0 & $\sim$8 GB & Chat general y RAG \\
LLaMA 3.1 8B Instruct & Q4\_0 & $\sim$4 GB & Resúmenes rápidos \\
Qwen 2.5 (LoRA FT)    & --- & $\sim$5 GB & Respuestas corporativas \\
Qwen 2.5 VL 7B        & --- & $\sim$5 GB & Análisis de imágenes \\
LLaVA 13B             & --- & $\sim$8 GB & OCR visual (legacy) \\
MiniLM-L12-v2          & FP32 & CPU only & Embeddings 384d \\
PaddleOCR              & --- & $\sim$500 MB & Reconocimiento de texto \\
\bottomrule
\end{tabular}
\end{table}

% ============================================================
\section{Resultados y Evaluación}
% ============================================================

\subsection{Métricas de Rendimiento}
El sistema fue sometido a pruebas de rendimiento sostenidas en un entorno de producción real. La Tabla~\ref{tab:metricas} sintetiza las métricas principales recogidas por Prometheus durante un período de operación continuada.

\begin{table}[H]
\centering
\caption{Métricas de rendimiento del sistema en producción.}
\label{tab:metricas}
\begin{tabular}{lc}
\toprule
\textbf{Métrica} & \textbf{Valor} \\
\midrule
Tiempo hasta primer token (TTFT) & $< 1.2$~s \\
Latencia RAG completa (p95)       & $\sim 3.5$~s \\
Throughput máximo sostenido       & $\sim 42$ req/min \\
Uptime del sistema                & 99.8\% (30 días) \\
Consumo VRAM pico (3 modelos)     & 18.5 / 24 GB \\
Alucinaciones en modo RAG         & 0\% (temp $\leq$ 0.2) \\
Cache hit rate (Redis)            & $\sim 35$\% \\
\bottomrule
\end{tabular}
\end{table}

\subsection{Evaluación Cualitativa}
Se realizaron pruebas funcionales abarcando los siguientes escenarios operativos:

\begin{enumerate}
    \item \textbf{Consulta RAG sobre documentos corporativos:} El usuario pregunta ``¿Cuál es la política de calidad de la empresa?'' y el sistema recupera los fragmentos más relevantes del PDF de calidad indexado, generando una respuesta fundamentada exclusivamente en el contenido real del documento.
    \item \textbf{Scraping de URL:} El usuario envía una URL y el sistema la descarga, renderiza con Playwright, limpia con Trafilatura, la trocea en chunks y genera una respuesta en contexto inmediato, almacenando además el contenido para posteriores consultas.
    \item \textbf{Búsqueda legal en el BOE:} El usuario solicita ``Busca en el BOE la ley de protección de datos'' y el agente resuelve semánticamente la ley, consulta la API oficial y devuelve un resumen estructurado.
    \item \textbf{Análisis de imagen:} El usuario adjunta una fotografía de un documento en papel y el sistema la procesa con los modelos de visión (Qwen 2.5 VL o LLaVA) para describir o transcribir su contenido.
\end{enumerate}

\subsection{Mitigación de Alucinaciones}
La eliminación de alucinaciones en modo RAG se logra mediante la combinación de dos técnicas:
\begin{itemize}
    \item \textbf{Temperatura determinista:} Los parámetros de inferencia se configuran con valores de temperatura extremadamente bajos (\texttt{0.1}--\texttt{0.2}), minimizando la aleatoriedad en la selección de tokens.
    \item \textbf{Prompt de sistema restrictivo:} El \textit{system prompt} inyectado instruye explícitamente al modelo para que responda \textbf{únicamente} basándose en el contexto recuperado, prohibiendo la generación libre de información no respaldada por los fragmentos.
\end{itemize}

El solapamiento de ventanas contextuales (\textit{overlap}) del 10--15\% en el chunking complementa estas medidas, asegurando que el hilo semántico de párrafos extensos no se pierda en los límites de fragmento.
